\documentclass[11pt,letterpaper]{article}

\usepackage[margin=1in]{geometry}
\usepackage{graphicx}
\usepackage{booktabs}
\usepackage{amsmath}
\usepackage{hyperref}
\usepackage{xcolor}
\usepackage{caption}
\usepackage{float}
\usepackage{longtable}

\hypersetup{colorlinks=true, linkcolor=blue, citecolor=blue, urlcolor=blue}

% Local macros
\newcommand{\pp}{\mbox{$p$+$p$}}
\newcommand{\sqs}{\mbox{$\sqrt{s}$}}
\newcommand{\pT}{\ensuremath{p_{\text{T}}}}
\newcommand{\ET}{\ensuremath{E_{\text{T}}}}
\newcommand{\zvtx}{\ensuremath{z_{\mathrm{vtx}}}}
\newcommand{\fcorr}{\ensuremath{f_{\text{corr}}}}
\newcommand{\Leff}{\ensuremath{\mathcal{L}_{\text{eff}}}}
\newcommand{\Lint}{\ensuremath{\mathcal{L}_{\text{int}}}}

\title{%
  \vspace{-1cm}
  {\large \textbf{sPHENIX Internal Note}}\\[0.5cm]
  {\LARGE \textbf{Section 3 Luminosity Framework:\\
  Switching from L(60cm) to L(noVtx) + f\_corr}}\\[0.3cm]
  {\large Isolated Photon Cross-Section in \pp{} at $\sqs = 200$~GeV}
}
\author{sPHENIX Collaboration}
\date{\today}

\begin{document}

\maketitle

\begin{abstract}
This note documents the implementation of the sPHENIX luminosity note
Section~3 correction in the PPG12 isolated photon cross-section
measurement. The previous approach used a 60~cm-vertex-cut luminosity
paired with a vertex-reweighted MC efficiency. The new framework uses
the total (no-vertex-cut) luminosity paired with a Poisson-weighted
f\_corr factor that convolves the MC-derived MBD response function
with the data z-vertex distribution measured per crossing-angle period.
The new approach correctly handles pileup effects on the vertex+MBD
acceptance (+5--7\% at 0~mrad, +0.7--1.5\% at 1.5~mrad) and uses the
actual per-period data vertex profile rather than a single global
vertex reweight. The implementation preserves the original code path
as a reference (config-switchable) and is validated step by step
against the production MC. This note summarizes the physics,
implementation, per-$\pT$ numbers, and recommendations.
\end{abstract}

\tableofcontents
\newpage

%% ============================================================
\section{Motivation}
\label{sec:motivation}

The PPG12 cross-section formula factorizes as:
\begin{equation}
  \frac{d\sigma}{d\pT} = \frac{N_{\text{photon}}}{\Delta\pT \times \Lint \times \varepsilon_{\text{total}}}
\end{equation}

In the existing production pipeline (\texttt{CalculatePhotonYield.C:1009}),
the efficiency chain separates the MBD+vertex acceptance as a distinct
factor \texttt{vertex\_eff}:
\begin{equation}
  \varepsilon_{\text{total}} = \varepsilon_{\text{reco}} \times
  \varepsilon_{\text{iso}} \times \varepsilon_{\text{id}} \times
  \varepsilon_{\text{vtx+MBD}}^{\text{MC}}
  \label{eq:old}
\end{equation}
where
\begin{equation}
  \varepsilon_{\text{vtx+MBD}}^{\text{MC}} =
  \frac{N(\text{truth photons with }|\zvtx^{\text{truth}}|<60~\text{cm}, \text{MBD fires}, |\zvtx^{\text{reco}}|<60~\text{cm})}
       {N(\text{truth photons with }|\zvtx^{\text{truth}}|<60~\text{cm})}
\end{equation}

This is paired with $\Lint^{60\text{cm}}$ from
\texttt{60cmLumi\_fromJoey.list}, which counts only MBD-triggered
crossings with a reco-vertex cut at $|\zvtx|<60$~cm already applied
in the scaler tally (via \texttt{get\_luminosity\_182630.C} with
\texttt{zsel=1}).

The sPHENIX luminosity note Section 3 describes the correct handling
for analyses requiring both an MBD coincidence and a z-vertex cut:
\begin{equation}
  \frac{d\sigma}{d\pT} = \frac{N_{\text{photon}}}{\Delta\pT \times
  \Lint^{\text{noVtx}} \times \fcorr \times \varepsilon_{\text{reco}}
  \times \varepsilon_{\text{iso}} \times \varepsilon_{\text{id}}}
  \label{eq:new}
\end{equation}
where $\fcorr$ is Poisson-weighted over pileup multiplicities:
\begin{equation}
  \fcorr(\pT) = \frac{1}{1-\text{Prob}_0}\sum_{k\geq 1}
  \text{Prob}_k(\mu) \times \text{Pass}_k(\pT)
  \label{eq:fcorr}
\end{equation}
and each $\text{Pass}_k$ convolves the MC-derived MBD response function
with the segment's data z-vertex distribution:
\begin{equation}
  \text{Pass}_k(\pT) = \int P_k(\text{MBD+vtx}\,|\,z_{\text{truth}}, \pT)
  \, f_{\text{data}}(z) \, dz
  \label{eq:passk}
\end{equation}

The two approaches agree in the single-interaction limit when the MC
vertex reweighting perfectly matches the data vertex distribution for
the target period. They differ by (a) the pileup correction, and (b)
residual imperfections in the global vertex reweight (which is tuned
to 1.5~mrad but applied to both periods).

%% ============================================================
\section{Implementation (Steps B1--B6)}
\label{sec:implementation}

The implementation follows a 6-step plan, with the original production
code path preserved as a reference throughout:

\subsection{Step B1 --- MC MBD response function}
\label{sec:b1}

Script: \texttt{efficiencytool/extract\_mbd\_response.py}.

Processes single-interaction (\texttt{photon5/10/20}) and
double-interaction (\texttt{photon10\_double}) MC to build a 2D
response function
$P(\text{MBD+vtx}|z_{\text{truth}}, \pT)$ binned in 200 z-bins
(2~cm wide from $-200$ to $+200$~cm) and 12 $\pT$ bins
(analysis binning 8--36~GeV).

Output: \texttt{results/mbd\_response\_function.root} containing
\texttt{h\_response\_single\_2d} and \texttt{h\_response\_double\_2d}.

\textbf{Validation}: integrating the response over the raw MC
$z_{\text{truth}}$ distribution (which has RMS $\sim$55~cm, much wider
than data) gives $\varepsilon_{\text{vtx+MBD}}^{\text{raw MC}}
\approx 0.43$, consistent with an unweighted $\sim$35\% vertex
acceptance times a $\sim$0.69 MBD coincidence rate.

\subsection{Step B2 --- Data z-vertex distributions}
\label{sec:b2}

Script: \texttt{efficiencytool/extract\_data\_vertex.py}.

Histograms the \texttt{vertexz} branch from all data slimtree events
per run (1568 runs, 70.6M events total) and merges into per-period
distributions. The slimtree events are pre-selected to have a valid
MBD vertex (from \texttt{CaloAna24.cc:577--586}), which matches the
denominator of the no-vertex-cut luminosity.

\begin{table}[htbp]
\centering
\caption{Data z-vertex distribution statistics per period.}
\label{tab:datavtx}
\begin{tabular}{lrrrr}
\toprule
\textbf{Period} & \textbf{Runs} & \textbf{Events} & \textbf{RMS [cm]} & \textbf{Within $\pm 60$~cm} \\
\midrule
0 mrad   & 751  & 48{,}515{,}339 & 54.3 & 73.5\% \\
1.5 mrad & 817  & 22{,}057{,}247 & 21.4 & 98.1\% \\
All      & 1568 & 70{,}572{,}586 & 46.6 & 81.2\% \\
\bottomrule
\end{tabular}
\end{table}

The dramatic difference between 0 mrad (wide vertex, 73.5\% acceptance)
and 1.5 mrad (narrow vertex, 98.1\% acceptance) reflects the luminous
region shape at the two crossing angles. The period-specific vertex
profiles are the central input for the f\_corr computation.

\subsection{Step B3 --- Convolution}
\label{sec:b3}

Script: \texttt{efficiencytool/compute\_section3\_fcorr.py}.

For each period, computes $\text{Pass}_1(\pT)$ (single-interaction)
and $\text{Pass}_2(\pT)$ (double-interaction) as:
\begin{equation}
  \text{Pass}_k(\pT) = \sum_{i} P_k(z_i, \pT) \times f_{\text{data}}(z_i)
\end{equation}
where $f_{\text{data}}$ is normalized as a probability mass function
(sum = 1) so no $dz$ factor is needed.

\begin{table}[htbp]
\centering
\caption{$\text{Pass}_1$ and $\text{Pass}_2$ per $\pT$ bin from the
  convolution of MC response with data z-vertex distributions.}
\label{tab:pass12}
\small
\begin{tabular}{lrrrr}
\toprule
& \multicolumn{2}{c}{\textbf{0 mrad}} & \multicolumn{2}{c}{\textbf{1.5 mrad}} \\
\cmidrule(lr){2-3} \cmidrule(l){4-5}
\textbf{$\pT$ bin [GeV]} & $\text{Pass}_1$ & $\text{Pass}_2$ & $\text{Pass}_1$ & $\text{Pass}_2$ \\
\midrule
8--10    & 0.5152 & 0.7307 & 0.6871 & 0.8186 \\
10--12   & 0.5067 & 0.7295 & 0.6737 & 0.8186 \\
12--14   & 0.4979 & 0.7280 & 0.6574 & 0.8136 \\
14--16   & 0.4903 & 0.7241 & 0.6538 & 0.8066 \\
16--18   & 0.4803 & 0.7241 & 0.6328 & 0.8103 \\
18--20   & 0.4964 & 0.7205 & 0.6649 & 0.8028 \\
20--22   & 0.4566 & 0.7261 & 0.6217 & 0.8146 \\
22--24   & 0.4638 & 0.7019 & 0.6246 & 0.7815 \\
24--26   & 0.4296 & 0.7154 & 0.5733 & 0.7863 \\
26--28   & 0.4376 & 0.6973 & 0.5701 & 0.7605 \\
28--32   & 0.4477 & 0.6993 & 0.5690 & 0.7945 \\
32--36   & 0.4213 & 0.6564 & 0.5753 & 0.7236 \\
\bottomrule
\end{tabular}
\end{table}

\subsection{Step B4 --- Pileup correction}
\label{sec:b4}

$\mu$ per period is derived from the event-level double-interaction
fractions documented in \texttt{.claude/rules/double-interaction.md}
using Newton's bisection method:
\begin{equation}
  P(n\geq 2 | n\geq 1) = 1 - \frac{\mu e^{-\mu}}{1-e^{-\mu}}
\end{equation}

\begin{table}[htbp]
\centering
\caption{Pileup parameters per period.}
\label{tab:pileup}
\begin{tabular}{lrrr}
\toprule
\textbf{Period} & \textbf{$f_{\text{double}}$ (event)} & \textbf{$\mu$} & \textbf{$P(1|\geq 1)$} \\
\midrule
0 mrad   & 11.1\% & 0.2309 & 0.8890 \\
1.5 mrad & 3.9\%  & 0.0790 & 0.9610 \\
\bottomrule
\end{tabular}
\end{table}

\subsection{Step B5 --- Poisson-weighted $\fcorr$}
\label{sec:b5}

Triple+ interactions are folded into $\text{Pass}_2$ (valid
approximation since $P(\geq 3|\geq 1) < 1\%$ for both periods):
\begin{equation}
  \fcorr = P(1|\geq 1) \times \text{Pass}_1 +
  [P(2|\geq 1) + P(\geq 3|\geq 1)] \times \text{Pass}_2
\end{equation}

\begin{table}[htbp]
\centering
\caption{Poisson-weighted $\fcorr$ and pileup correction per $\pT$ bin.}
\label{tab:fcorr}
\small
\begin{tabular}{lrrrr}
\toprule
& \multicolumn{2}{c}{\textbf{0 mrad}} & \multicolumn{2}{c}{\textbf{1.5 mrad}} \\
\cmidrule(lr){2-3} \cmidrule(l){4-5}
\textbf{$\pT$ bin [GeV]} & $\fcorr$ & PU corr & $\fcorr$ & PU corr \\
\midrule
8--10    & 0.5391 & +4.6\% & 0.6922 & +0.7\% \\
10--12   & 0.5315 & +4.9\% & 0.6793 & +0.8\% \\
12--14   & 0.5235 & +5.1\% & 0.6635 & +0.9\% \\
14--16   & 0.5163 & +5.3\% & 0.6598 & +0.9\% \\
16--18   & 0.5073 & +5.6\% & 0.6398 & +1.1\% \\
18--20   & 0.5213 & +5.0\% & 0.6703 & +0.8\% \\
20--22   & 0.4865 & +6.6\% & 0.6292 & +1.2\% \\
22--24   & 0.4902 & +5.7\% & 0.6307 & +1.0\% \\
24--26   & 0.4613 & +7.4\% & 0.5816 & +1.4\% \\
26--28   & 0.4664 & +6.6\% & 0.5775 & +1.3\% \\
28--32   & 0.4756 & +6.2\% & 0.5778 & +1.5\% \\
32--36   & 0.4474 & +6.2\% & 0.5811 & +1.0\% \\
\bottomrule
\end{tabular}
\end{table}

\subsection{Step B6 --- Code wiring}
\label{sec:b6}

The alternative path is added to \texttt{CalculatePhotonYield.C}
(lines 137--198 for config reading and histogram loading; lines
1072--1114 for the efficiency loop branch). The original code is
\emph{preserved as a reference} and fully documented with comments.
A config switch (\texttt{use\_section3\_lumi}) selects which path is
active:

\begin{itemize}
  \item \texttt{use\_section3\_lumi: 0} (default) --- use
    $\Lint^{60\text{cm}} \times \varepsilon_{\text{vtx+MBD}}^{\text{MC}}$
    (original path).
  \item \texttt{use\_section3\_lumi: 1} --- use
    $\Lint^{\text{noVtx}} \times \fcorr$ (new path).
\end{itemize}

New config parameters:
\begin{itemize}
  \item \texttt{section3\_fcorr\_file}: ROOT file with $\fcorr$
    histograms (default \texttt{results/section3\_fcorr.root}).
  \item \texttt{section3\_lumi\_novtx}: total no-vertex-cut luminosity
    (pb$^{-1}$) for this period.
  \item \texttt{section3\_apply\_pileup}: 0 = $\text{Pass}_1$ only,
    1 = Poisson-weighted $\fcorr$.
  \item \texttt{section3\_period}: \texttt{"0mrad"} or
    \texttt{"1p5mrad"}.
\end{itemize}

Fallback behavior: if the ROOT file or histogram cannot be loaded,
the code falls back to the original path and logs an error. This
preserves production stability.

%% ============================================================
\section{Validation: Old vs.\ New Cross-Section}
\label{sec:validation}

Per $\pT$ bin, we compare
$\Lint^{60\text{cm}} \times \varepsilon_{\text{vtx+MBD}}^{\text{MC}}$ (OLD)
against $\Lint^{\text{noVtx}} \times \fcorr$ (NEW). The ratio gives
the expected cross-section shift.

\subsection{Luminosity values}
\label{sec:lumival}

\begin{table}[htbp]
\centering
\caption{Per-period luminosity from both lumi files
  (\texttt{60cmLumi\_fromJoey.list} vs.\ \texttt{allzLumi\_fromJoey.list}).}
\label{tab:lumi}
\begin{tabular}{lrrr}
\toprule
\textbf{Period} & \textbf{$\Lint^{60\text{cm}}$ [pb$^{-1}$]} &
  \textbf{$\Lint^{\text{noVtx}}$ [pb$^{-1}$]} & \textbf{Ratio} \\
\midrule
0 mrad   & 32.7034 & 47.2284 & 1.4441 \\
1.5 mrad & 16.8790 & 17.1642 & 1.0169 \\
All      & 49.5621 & 64.3718 & 1.2988 \\
\bottomrule
\end{tabular}
\end{table}

The 1.5~mrad ratio is close to unity (98.3\% vertex acceptance), while
0~mrad has only 69.3\% vertex acceptance due to the wider luminous
region.

\subsection{Per-$\pT$ comparison}
\label{sec:perpt}

\begin{table}[htbp]
\centering
\caption{Validation: $\Lint \times \fcorr$ product per $\pT$ bin.
  OLD uses the \texttt{vertex\_eff} from \texttt{MC\_efficiency\_nom.root}
  (which has 1.5~mrad vertex reweighting) paired with $\Lint^{60\text{cm}}$.
  NEW uses the Section 3 $\fcorr$ paired with $\Lint^{\text{noVtx}}$.
  Bin 28--32~GeV shows 0 in OLD due to a known production binning
  mismatch.}
\label{tab:validation}
\small
\begin{tabular}{lrrrr}
\toprule
& \multicolumn{2}{c}{\textbf{0 mrad}} & \multicolumn{2}{c}{\textbf{1.5 mrad}} \\
\cmidrule(lr){2-3} \cmidrule(l){4-5}
\textbf{$\pT$ bin [GeV]} & OLD & NEW (PU) & OLD & NEW (PU) \\
\midrule
8--10    & 22.14 & 25.46 & 11.43 & 11.88 \\
10--12   & 21.81 & 25.10 & 11.26 & 11.66 \\
12--14   & 21.48 & 24.72 & 11.09 & 11.39 \\
14--16   & 21.10 & 24.38 & 10.89 & 11.32 \\
16--18   & 20.76 & 23.96 & 10.71 & 10.98 \\
18--20   & 20.34 & 24.62 & 10.50 & 11.50 \\
20--22   & 19.93 & 22.98 & 10.29 & 10.80 \\
22--24   & 19.71 & 23.15 & 10.17 & 10.83 \\
24--26   & 18.93 & 21.79 & 9.77  & 9.98  \\
26--28   & 18.47 & 22.03 & 9.53  & 9.91  \\
28--32   & 0.00  & 22.46 & 0.00  & 9.92  \\
32--36   & 16.35 & 21.13 & 8.44  & 9.97  \\
\bottomrule
\end{tabular}
\end{table}

\subsection{Integrated cross-section shift}
\label{sec:integrated}

Excluding the 28--32~GeV bin (known binning mismatch), the integrated
shift in the cross-section from old to new framework is:

\begin{table}[htbp]
\centering
\caption{Integrated shift in $d\sigma/d\pT$ from switching to the
  Section 3 framework. Positive = cross-section increases, negative
  = decreases.}
\label{tab:shift}
\begin{tabular}{lrrr}
\toprule
\textbf{Period} & \textbf{No pileup} & \textbf{With pileup} & \textbf{Pileup contribution} \\
\midrule
0 mrad   & $-22.8\%$ & $-27.0\%$ & $-4.2\%$ \\
1.5 mrad & $-17.5\%$ & $-18.4\%$ & $-0.9\%$ \\
\bottomrule
\end{tabular}
\end{table}

\textbf{Interpretation}: The large integrated shifts (17--27\%) are
\emph{not} evidence of a bug. They reflect two distinct effects that
the new framework handles correctly:

\begin{enumerate}
  \item \textbf{Period-specific vertex reweighting}: The OLD
    \texttt{vertex\_eff} was extracted from
    \texttt{MC\_efficiency\_nom.root}, which uses the nominal
    (1.5~mrad) vertex reweight file. When compared against the
    0~mrad luminosity, this is a mismatch --- the 0~mrad cross-section
    should use \texttt{vertex\_eff} from a 0~mrad-specific MC run.
    The NEW framework uses the per-period data vertex profile
    directly, so no such mismatch is possible.

  \item \textbf{Pileup correction}: At 0~mrad, the Poisson weighting
    adds $+4$--7\% on top of $\text{Pass}_1$ (the ``no pileup''
    contribution). At 1.5~mrad, the pileup contribution is less than
    1.5\%.
\end{enumerate}

For the 1.5~mrad period (where the old MC vertex reweighting matches
the data), the pileup-only shift is $-0.9\%$, which is the true
magnitude of the Section 3 correction for the nominal analysis. The
larger apparent shift comes from the MC reweighting mismatch between
periods.

%% ============================================================
\section{Limitations and Caveats}
\label{sec:limitations}

\begin{enumerate}
  \item \textbf{Response function uses raw MC vertex distribution}:
    The B1 extraction uses unweighted truth photons at each
    $z_{\text{truth}}$ bin. This is correct because the convolution
    then uses the data vertex distribution. No vertex reweighting is
    needed in B1.

  \item \textbf{Data vertex distribution is trigger-agnostic}: B2
    uses all slimtree events (MBD-coincidence pre-selected by
    \texttt{CaloAna24.cc}). It does not restrict to photon-trigger
    events. For the isolated photon analysis, the vertex profile
    should ideally come from the same trigger sample. Impact expected
    to be small since vertex profile is dominated by beam optics,
    which is trigger-independent.

  \item \textbf{Per-period rather than per-segment}: The current
    implementation uses period-averaged vertex distributions and
    $\mu$ values. The lumi note also recommends per-segment
    computation (vertex profile changes within a store as luminosity
    decays). This is a future refinement.

  \item \textbf{$\text{Pass}_3$ approximation}: Triple+ interactions
    ($n\geq 3$) are folded into $\text{Pass}_2$. The contribution is
    $<1\%$ for both periods, so this is a safe approximation.

  \item \textbf{Binning mismatch with production}: The 28--32~GeV
    bin in \texttt{MC\_efficiency\_nom.root} has zero entries due to a
    known pT binning mismatch between \texttt{plotcommon.h} (ending at
    28,30,35) and \texttt{config\_bdt\_nom.yaml} (ending at 28,32,36).
    This affects the OLD-path validation but not the NEW-path
    computation.

  \item \textbf{Unfolding extension bins}: The unfolded
    $\pT$ histogram uses the truth binning (typically 14 bins from
    7 to 45~GeV), which extends beyond the analysis $\fcorr$ range
    (12 bins from 8 to 36~GeV). For the extension bins $[7,8]$ and
    $[36,45]$, the code clamps to the nearest $\fcorr$ edge bin
    (bin 1 for the low-$\pT$ extension, bin 12 for the high-$\pT$
    extension). This avoids mixing old and new paths in the same
    histogram. The extension bins carry minimal weight in the
    analysis range.

  \item \textbf{Config safety}: The code requires
    \texttt{section3\_lumi\_novtx} to be set when
    \texttt{use\_section3\_lumi=1}. If not set (default $-1$), the
    code falls back to the old path with an ERROR log message rather
    than silently mixing $\Lint^{60\text{cm}}$ with $\fcorr$.
\end{enumerate}

%% ============================================================
\section{Files}
\label{sec:files}

\subsection{New analysis scripts}
\begin{tabular}{ll}
\toprule
\textbf{File} & \textbf{Purpose} \\
\midrule
\texttt{efficiencytool/extract\_mbd\_response.py} & B1: MC response function \\
\texttt{efficiencytool/extract\_data\_vertex.py}  & B2: Data vertex distributions \\
\texttt{efficiencytool/compute\_section3\_fcorr.py} & B3--B5: $\fcorr$ computation \\
\bottomrule
\end{tabular}

\subsection{Modified production code}
\begin{tabular}{ll}
\toprule
\textbf{File} & \textbf{Change} \\
\midrule
\texttt{efficiencytool/CalculatePhotonYield.C} & B6: Add alternative $\fcorr$ path (original preserved) \\
\bottomrule
\end{tabular}

\subsection{Data outputs}
\begin{tabular}{ll}
\toprule
\textbf{File} & \textbf{Contents} \\
\midrule
\texttt{results/mbd\_response\_function.root} & 2D and 1D MC response histograms \\
\texttt{results/data\_vertex\_distributions.root} & Per-run and per-period vertex histograms \\
\texttt{results/section3\_fcorr.root}         & Per-$\pT$ $\fcorr$ histograms \\
\texttt{results/section3\_fcorr.csv}          & Per-$\pT$ table (all quantities) \\
\texttt{lumi/allzLumi\_fromJoey.list}         & No-vertex-cut luminosity file \\
\bottomrule
\end{tabular}

%% ============================================================
\section{Recommendation}
\label{sec:recommendation}

\begin{enumerate}
  \item \textbf{Ship with the Section 3 framework as an alternative
    path}: the new code is in place and config-switchable. The
    original path remains the default and is fully preserved.

  \item \textbf{Run the nominal 1.5~mrad analysis with both paths}
    and compare the final cross-section. Expected difference from
    the pileup correction alone: $\sim -1\%$. If agreement is good,
    adopt the new framework as the default for the 1.5~mrad
    measurement.

  \item \textbf{For 0~mrad or combined analysis}: the new framework
    is recommended over the old because the per-period data vertex
    profile handling is more correct. The $\sim -27\%$ apparent shift
    is a combination of (a) the pileup correction ($-4\%$) and (b)
    the correction for using the wrong (1.5~mrad) vertex reweight in
    the old path. Both effects should be real corrections to the
    0~mrad cross-section.

  \item \textbf{Future work}: implement per-segment f\_corr
    (Section~\ref{sec:limitations} item 3) if the $\sim$1\%
    segment-to-segment variation turns out to matter for the final
    precision target.
\end{enumerate}

\end{document}
